\documentclass[../../main.tex]{subfiles}
\begin{document}

\begin{flushright}
\footnotesize\textit{【自動文字起こし・要確認】}
\end{flushright}

\noindent
[\textbf{解}] $0 < x < 1 \dots \text{①}$

\[
\dfrac{1}{x} \int_0^x \dfrac{1}{\sqrt{1-t}} \, dt = \dfrac{1}{\sqrt{1-y}} \quad \dots \text{②}
\]

\[
\dfrac{1}{x} \int_0^x \dfrac{1}{\sqrt{1-t}} \, dt = \dfrac{1}{x} \left[ -2(1-t)^{\dfrac{1}{2}} \right]_0^x = \dfrac{-2}{x} \left(\sqrt{1-x} - 1\right)
\]

を②に代入

\[
\dfrac{2}{x} \left(1 - \sqrt{1-x}\right) = \dfrac{1}{\sqrt{1-y}}
\]

①では両辺正だから2乗して

\[
\dfrac{4}{x^2} \left(1 - \sqrt{1-x}\right)^2 = \dfrac{1}{1-y}
\]

同じく逆数をとって

\[
\dfrac{x^2}{4\left(1 - \sqrt{1-x}\right)^2} = 1 - y
\]

\[
\therefore y = 1 - \dfrac{x^2}{4\left(1 - \sqrt{1-x}\right)^2} \equiv f(x) \quad \dots \text{③}
\]

さて,
\[
\begin{aligned}
4 f'(x) &= -\dfrac{2x\left(1-\sqrt{1-x}\right)^2 - x^2 \cdot 2\left(1-\sqrt{1-x}\right) \dfrac{-1}{2\sqrt{1-x}}}{\left(1-\sqrt{1-x}\right)^4} \\
&= \dfrac{-2x}{\left(1-\sqrt{1-x}\right)^3} \left[ \left(1-\sqrt{1-x}\right) - \dfrac{x}{2\sqrt{1-x}} \right] \\
&= \dfrac{-x}{\left(1-\sqrt{1-x}\right)^3 \sqrt{1-x}} \left[ 2\sqrt{1-x} - (2-x) \right] \\
&= \dfrac{1}{\left(1-\sqrt{1-x}\right)^3 \sqrt{1-x}} \dfrac{x^3}{2\sqrt{1-x} + (2-x)} > 0 \quad (\because 0 < x < 1)
\end{aligned}
\]

から, $y$ は $x$ の単調増加であり, 又, ③の表すから

\[
y \longrightarrow \dfrac{3}{4} \quad (x \to 1)
\]

である.

\end{document}
